\section{Operatory Hilberta-Schmidta}

Wyjątkowo w tych notatkach zacznijmy od historii. Na przełonie XX wieku Hilbert pracował na przestrzeniach z iloczynem skalarnym (później nazwanych przestrzeniami Hilberta) takimi jak $\ell^2(\N)$. Jest to przestrzeń ciągów $\{a_n\}_n\geq 0$ sumowalnych z kwadratem, czyli $\sum_{n\geq0}|a_n|^2<\infty$. Jest to przestrzeń nieskończenie wymiarowa, więc naturalna jest próba zdefiniowania nieskończenie wymiarowej wersji macierzy. Jego student, Schmidt, rozważał macierze $\{m_{i,j}\}$ takie, że $\sum |m_{j,k}|^2<\infty$. W ten sposób dostajemy pierwsze przykłady operatorów Hilberta-Schmidta.

Rozważamy unitarne reprezentacje grup zwartych $G$ z unormowaną miarą Haara $dx$, to znaczy $dx(G)=1$.

% \begin{definition}{H-*-algebra}{}
%   Algebra $A$ z inwolucją $x\mapsto x^*$ która jest przestrzenią Hilberta z iloczynem skalarnym $\langle \cdot,\cdot\rangle$ nazywa się \buff{H-*-algebrą} jeśli 
%   \begin{itemize}
%     \item $(\forall\;x,y,z\in A)\;\langle xy, z\rangle=\langle y, x^*z\rangle$
%     \item $\|x^*\|=\|x\|$
%   \end{itemize}
%   Dodatkowo chcemy, żeby ta algebra była \acc{półprosta}, czyli dla $x\neq 0$ wymagamy $xx^*\neq 0$.
%     %
%     % Powiemy, że przestrzeń Hilberta $A$ z unitarnym działaniem $*$ nazywamy *-algebrą jeśli
%     % \begin{itemize}
%     %   \item $(\forall\;x,y,z\in A)\;\langle xy, z\rangle=\langle y, x^*z\rangle$
%     %   \item $\|x^*\|=\|x\|$
%     % \end{itemize}
% \end{definition}
%
% \begin{definition}{algebra Hilberta}{}
%   W literaturze algebrą Hilberta nazywamy algebrę $A$ z inwolucją $x\mapsto x^*$ oraz iloczynem skalarnym $\langle \cdot,\cdot\rangle$ taką, że
%   \begin{itemize}
%     \item $(\forall\;x,y,z\in A)\;\langle xy, z\rangle=\langle y, x^*z\rangle$
%     \item $(\forall\;x,y\in A)\;\langle x,y\rangle=\langle y^*,x^*\rangle$
%     \item dla dowolnego ustalonego $a\in A$ operator $x\mapsto ax$ jest ograniczony
%     \item przestrzeń liniowa rozpięta przez $xy$ dla $x,y\in A$ jest gęstym podzbiorem $A$.
%   \end{itemize}
% \end{definition}

Niech $T:\Hh\to\Hh$ będzie operatorem ograniczonym na przestrzeni Hilberta $\Hh$ z bazą $\{e_i\;:\;i\in I\}$. Definiujemy normę Hilberta-Schmidta operatora $T$ jako
$$\|T\|_{HS}^2:=\sum_{i\in I}\|Te_i\|_H^2.$$
Powiemy, że $T$ jest \buff{operatorem Hilberta-Schmidta} jeśli ma skończoną normę Hilberta-Schmidta.

\begin{fact}{}{}
  Norma Hilbera-Schmidta jest dobrze określona, tzn. nie zależy od wyboru bazy przestrzeni $\Hh$.
\end{fact}

\begin{proof}
  Zaczynamy od pokazania, że $\|T\|_{HS}^2=\|T^*\|_{HS}^2$. Pitagoras mówi nam, że jeśli $\{e_i\;:\;i\in I\}$ jest bazą ortonormalną $\Hh$, to
  $$\|x\|^2=\sum_i|\langle x, e_i\rangle|^2=\sum_i|\langle e_i, x\rangle|^2,$$
  bo $x=\sum\langle x,e_i\rangle e_i$. Stąd
  $$\|T\|_{HS}^2=\sum_i\|Te_i\|^2_{HS}=\sum_i\sum_j|\langle Te_i, e_j\rangle|^2=\sum_i\sum_j|\langle e_i, T^*e_j\rangle|^2=\sum_j\|T^*e_j\|^2=\|T^*\|_{HS}^2.$$
  Niech teraz $\{f_i\;:\;i\in I\}$ będzie konkurencyjna bazą ortonormalną $\Hh$. Wówczas
  $$\sum_i\|Te_i\|^2=\sum_i\sum_j|\langle Te_i, f_j\rangle|^2=\sum_i\sum_j|\langle e_i, T^*f_j\rangle|^2=\sum_j\sum_i|\langle e_i, T^*f_j\rangle|^2=\sum_j\|T^*f_j\|^2.$$
\end{proof}

\begin{remark}{}{}
  Niech $B_{HS}(\Hh)$ będzie zbiorem wszystkich operatorów Hilberta-Schmidta na $\Hh$. Wówczas $B_{HS}(\Hh)$ jest obustronnym $*$-ideałem w $B(\Hh)$.
\end{remark}

Zauważmy, że
$$\|T\|_{HS}^2=\sum_{i\in I}\langle Te_i, Te_i\rangle=\tr(T*T),$$
odważnym posunięciem jest więc zdefiniowanie iloczynu Hilberta-Schmidta jako
$$\langle T,S\rangle_{HS}:=\tr(S^*T)=\sum_{i\in I}\langle Te_i, Se_i\rangle.$$

\begin{remark}{}{}
  Przestrzeń operatorów Hilberta-Schmidta z iloczynem $\langle\cdot,\cdot\rangle_{HS}$ tworzy przestrzeń Hilberta.
\end{remark}

\begin{theorem}{}{}
  $T$ jest operatorem Hilberta-Schmidta wtedy i tylko wtedy gdy $|T|=\sqrt{T^*T}$ jest operatorem Hilberta-Schmidta.
\end{theorem}

\begin{definition}{norma Schattena}{}
  Niech $p\in[1,\infty)$ oraz $T\in B(\Hh)$ będzie operatorem ograniczonym. Wówczas definiujemy \buff{$p$-normę Schattena} jako
  $$\|T\|_p:=[\tr(|T|^p)]^{\frac{1}{p}}.$$
\end{definition}

Dla $p=2$ norma Schattena jest normą Hilberta-Schmidta, dla $p=1$ dostajemy normę śladową (nuklearną), a dla $p=\infty$ dostajemy $\|T\|_{\infty}=\sup_{\|\xi\|=1}\|T\xi\|$ zwykłą normę operatorową.

\begin{lemma}{}{}
  $$\|T\|_{HS}\geq \|T\|$$
\end{lemma}

\begin{proof}
  Dla dowolnego $\epsilon>0$ niech $\xi_0$, $\|\xi_0\|=1$, będzie wektorem takim, że 
  $$\|T\xi_0\|\geq \|T\|-\epsilon.$$
  Możemy takie $\xi_0$ rozszerzyć do bazy ortonormalnej $\{\xi_i\;:\;i\in I\}$ by dostać
  $$\|T\|_{HS}^2=\sum_i\|T\xi_i\|^2\geq\|T\xi_0\|^2\geq (\|T\|-\epsilon)^2.$$
  Przechodząc z $\epsilon\to 0$ dostajemy pożądaną nierówność.
\end{proof}

\begin{example}[m]
  \item Rozważmy $\Hh=L^2(X,\mu)$ z $\mu(X)<\infty$. Dla $k\in L^2(X\times X)$ rozważmy operator
    $$(T_kf)(x)=\int_Xk(x,y)f(y)d\mu(y).$$
    Jego norma Hilberta-Schmidta zgadza się z $L^2$ normą $k$ (ćwiczenie):
    $$\|T_k\|^2_{HS}=\int_X\int_X |k(x,y)|^2 d\mu(y)d\mu(x).$$
    \bigskip 

    Co więcej, jeśli $k_1, k_2\in L^2(X\times X)$, to zachodzi
    $$T_{k_1}T_{k_2}=T_{k_3},$$
    gdzie
    $$k_3(x,y)=\int_X k_1(x,z)k_2(z,y)d\mu(z).$$
  \item Dla $\phi\in C(G)$ określamy operator splotu 
    $$(T_\phi f)(x)=(f\ast\phi)(x)=\int_G f(y)\phi(y^{-1}x)dy$$
\end{example}


\begin{theorem}{}{}\mbox{}
  \begin{enumerate}
    \item Operatory o normie Schattena $1$ tworzą $*$-ideał w $B(\Hh)$.
    \item Co więcej, tworzą one przestrzeń Banacha.
  \end{enumerate}
\end{theorem}

\subsection{O zbieżności słabych słów kilka}

\begin{definition}{słaba zbieżność}{}
  Powiemy, że ciąg $x_n$ z przestrzeni Hilberta $\Hh$ zbiega słabo do $x$, jeśli 
  $$(\forall\;y\in\Hh)\langle x_n, y\rangle\to\langle x,y\rangle$$
\end{definition}

Przez mocną zbieżność będziemy rozumieć zbieżność w normie, tzn. $x_n$ zbiega mocno do $x$ jeśli $\|x_n-x\|\to 0$.

\begin{remark}{}{}
  Jeśli $x_n$ zbiega słabo do $x$ oraz $\|x_n\|\to \|x\|$ to $x_n$ zbiega mocno do $x$.
\end{remark}

\begin{proof}
  \begin{align*}
    \|x_n-x\|^2&=\langle x_n-x, x_n-x\rangle=\\ 
               &=\langle x_n, x_n\rangle -2Re\langle x_n, x\rangle + \langle x, x\rangle=\\ 
               &=\|x_n\|^2-2Re\langle x_n, x\rangle +\|x\|^2\to\\ 
               ^\to \|x\|^2-2Re\langle x, x\rangle + \|x\|^2=\\ 
               &=2\|x\|^2-2\|x\|^2=0
  \end{align*}
\end{proof}

\begin{theorem}{nierówność Bessel'a}{}
  Jeśli $\{e_n\}$ to baza ortonormalna $\Hh$, to dla dowolnego wektora $x\in\Hh$ zachodzi
  $$\sum_{k\geq 0}|\langle x, e_k\rangle|^2\leq \|x\|^2$$
\end{theorem}

\begin{proof}
  Dla dowolnego $N$ niech 
  $$x_N=\sum_{k\leq N}\langle x, e_k\rangle e_k,$$
  co jest rzutem ortogonalnym $x$ na podprzestrzeń $\langle e_0,...,e_k\rangle$, czyli $x-x_N\perp x_N$. Twierdzenie Pitagorasa mówi nam, że $\|x\|^2=\|x_N\|^2+\|x-x_N\|^2$, czyli wystarczy wyliczyć $\|x-x_N\|$:
  \begin{align*}
    0&\leq \|x-x_N\|^2=\left\langle x-\sum_{k\leq N}\langle x, e_k\rangle e_k, x-\sum_{j\leq N}\langle x, e_j\rangle e_j\right\rangle=\\ 
                     &=\|x\|^2-\sum_{k\leq N}\overline{\langle x, e_k\rangle}\langle x, e_k\rangle-\sum_{j\leq N}\langle x, e_j\rangle \langle x, e_j\rangle+\sum_{j\leq N}\sum_{k\leq N}\overline{\langle x, e_k\rangle}\langle x, e_j\rangle\langle e_k, e_j\rangle=\\ 
                     &=\|x\|^2-\sum_{k\leq N}|\langle x, e_k\rangle|^2
  \end{align*}
  Dostajemy $\|x\|^2\geq \sum_{k\leq N}|\langle x, e_k\rangle|^2$, czyli w granicy również ta nierówność zachodzi.
\end{proof}

\begin{lemma}{}{slaba zbieznosc iff po wspolrzednych}
  Niech $\{e_n\}$ będzie baza ortonormalną przestrzeni $\Hh$. Wówczas ograniczony ciąg $x_n$ zbiega słabo do $x$ wtedy i tylko wtedy gdy
  $$(\forall\;m)\;\langle x_n, e_m\rangle \to \langle x, e_m\rangle.$$
\end{lemma}

\begin{proof}
  $\implies$ 

  $e_n$ jest dowolnym elementem $\Hh$. 

  $\impliedby$

  Niech $M\geq\|x_n\|$ dla wszystkich $n$. Ustalmy dowolne $y\in \Hh$ i zapiszmy je w bazie $\{e_n\}$ jako $y=\sum y_ne_n$. Ponieważ $\|y\|<\infty$ to dla każdego $\epsilon>0$ istnieje $N$ takie, że 
  $$\left \|\sum_{n\geq N} y_n e_n \right \|<\epsilon.$$
  Z założenia dostajemy
  $$\lim_m \langle x_m, \sum_{n\leq N} y_ne_n\rangle=\sum_{n\leq N}y_n\lim_m \langle x_m, e_n\rangle=\sum_{n\leq N}y_n\langle x, e_n\rangle=\langle x, \sum_{n\leq N}y_ne_n\rangle.$$
  Przechodząc do sedna,
  \begin{align*}
    &|\langle x_m, y\rangle-\langle x, y\rangle|=\\ 
    =&|\langle x_m, y\rangle-\langle x_m, \sum_{n\leq N}y_ne_n\rangle+\langle x_m, \sum_{n\leq N}y_ne_n\rangle-\langle x, \sum_{n\leq N}y_ne_n\rangle+\langle x, \sum_{n\leq N}y_ne_n\rangle-\langle x, y\rangle|\leq \\ 
    \leq& |\langle x_m, \sum_{n> N}y_ne_n\rangle|+|\langle x_m-x, \sum_{n\leq N}y_ne_n\rangle|+|\langle x, \sum_{n>N}y_ne_n\rangle|
  \end{align*}
  środkowy wyraz zbiega do $0$, a skrajne dzięki nierówności Cauchy'ego-Schwarza są dowolnie małe. Stąd $\langle x_m, y\rangle\to\langle x, y\rangle$.
\end{proof}

Dla przypomnienia szybki lemat z analizy funkcjonalnej.

\begin{lemma}{}{kula jednostkowa slabo zwarta}
  Kula jednostkowa $B$ jest zwarta w słabej topologii.
\end{lemma}

\begin{proof}
  Niech $x_n\in B$ będzie dowolnym ciągiem, a $\{e_n\}$ będzie baza ortonormalną przestrzeni Hilberta w której pracujemy. Ponieważ liczby $\langle e_1, x_m\rangle$ tworzą ciąg z $[-1, 1]$ to możemy wybrać podciąg $x_{\alpha_1(m)}$ dla którego $\langle e_1, x_{\alpha_1(m)}\rangle$ zbiega do $a_1$. Powtarzamy to samo dla $\langle e_2, x_{\alpha_1(m)}\rangle$ i dostajemy podciąg $x_{\alpha_2(m)}$ dla którego te liczby zbiegają do $a_2$. Iterując tę procedurę dostajemy ciąg ciągów $x_{\alpha_n(m)}$ oraz liczb $a_n$ do których one zbiegają na $n$-tej współrzędnej. 

  Interesuje nas podciąg z przekątnej tej konstrukcji, tzn. 
  $$y_m=x_{\alpha_m(m)}.$$
  Zauważmy, że dla każdego $N$ zachodzi
  $$\sum_{n\leq N}|a_n|^2=\lim_m\sum_{n\leq N}|\langle e_n,y_m\rangle|^2\leq \lim\|y_\|^2\leq 1.$$
  Stąd $\sum |a_n|^2\leq 1$ i możemy zdefiniować kandydata na granicę podciągu $y_m$ ciągu $x_n$ w następujący sposób
  $$y=\sum a_n e_n\in B.$$
  Ponieważ $y_n\to y$ po współrzędnych, to z lematu \ref{lmm:slaba zbieznosc iff po wspolrzednych}.
\end{proof}

\begin{theorem}{}{operatory hilberta schmidta sa calkowe}
    Każdy operator Hilberta-Schmidta jest postaci $Af(x)=\int_X k(x,y)f(y)d\mu(y)$ dla pewnego $k\in L^2(X\times X)$, to znaczy
    $$\langle f\;|\;Ag\rangle=\int\int \vec{f}(x)k(x,y)g(y)dxdy.$$
\end{theorem}


\begin{proof}
  Zaczynamy od powiedzenia, że każdy ograniczony operator $A$ którego obraz ma skończony wymiar zapisuje się jako $Af(x)=\int k(x,y)f(y)dy$. Dowód tego stwierdzenia pozostawiamy jako ćwiczenie, które powinno się zacząć od tego, że $\dim(\im(A))<\infty\iff\dim(\im(A^*))<\infty$ oraz zdefiniowania $k(x,y)=\sum_jA\phi_j(x)\phi_j(y)$ dla $\{\phi_i\}$ bazy $\im(A^*)$.

    Niech $A$ będzie teraz operatorem Hilberta-Schmidta, natomiast $A_n$ niech będzie ciągiem operatorów o skończenie wymiarowym obrazie takich, że $\|A_n\|\to \|A\|$. Odwzorowanie $k\mapsto T_k$, gdzie $T_k$ jest operatorem w postaci przez nas pożądanej, jest izometrią (ćwiczenie). Niech $T_{k_n}=A_n$, wtedy $k_n$ będzie ciągiem Cauchy'ego, więc jest przenoszony na ciąg Cauchy'ego $T_{k_n}$, który musi mieć granicę zgadzającą się z granicą $A_n$. Stąd $k=\lim k_n$ daje nam tezę.
\end{proof}

\begin{lemma}{}{a jest ciagla}
  Dla każdego operatora Hilberta-Schmidta funkcjonał na kuli jednostkowej $B\ni f\mapsto \langle Af, f\rangle\in\C$ jest ciągły.
\end{lemma}

Funkcję $f\mapsto \langle Af, f\rangle$ oznaczymy $a(f)$.

\begin{proof}
  \color{red} Jeśli $\{e_n\}$ jest bazą ortonormalną w $L^2X$, to $\{e_ne_m\}$ tworzy bazę ortonormalną $L^2(X\times x)$. W takim razie
    $$k(x,y)=\sum a_{n,m}e_n(x)e_m(y),$$
    gdzie $\sum|a_{n,m}|^2<\infty$.
\end{proof}

\subsection{Twierdzenie spektralne dla Hilberta-Schmidta}

Celem tego rozdzialiku będzie przedstawienie, wraz z dowodem, twierdzenia spektralnego dla samosprzeżonych operatorów Hilberta-Schmidta. Zanim do tego przejdziemy jednak potrzebujemy szybki lemacik.

\begin{lemma}{}{supremum wo wartosc wlasna}
  Niech $A=A^*$ będzie operatorem samosprzężonym. Jeśli $x_0\in B$ taki, że
  $$\langle Ax_0, x_0\rangle=a(x_0)=\sup_{y\in B}a(y),$$
  to $\|x_0\|=1$ i jest on wektorem własnym $A$, tzn. istnieje $\lambda\in\C$ takie, że $Ax_0=\lambda x_0$.
\end{lemma}

Analogicznie zachowuje się $x_0$ takie, że $a(x_0)=\inf_{y\in B}a(y)$.

\begin{proof}
  Zacznijmy od udowodnienia, że $\|x_0\|=1$. Dla dowolnego $x\in B$ zachodzi
  $$a(x)=\|x\|^2a(\frac{x}{\|x\|}).$$
  W takim razie jeśli $\|x\|<1$ to $a(x)<a(\frac{x}{\|x\|})$, czyli nie osiągnie on nigdy supremum. Stąd $\|x_0\|=1$.

  Ustalmy $y\in B$ takie, że $\|y\|=1$ oraz $\langle x_0,y\rangle=0$. Wtedy dla każdego $\alpha\in\R$ zachodzi $\|x_0+\alpha y\|^2=1+\alpha^2$. Skoro w $x_0$ mamy maksimum funkcjonału $a$, to pochodna musi się zerować 
  \begin{align*}
    0&=\frac{\dd}{\dd\alpha}\left[a\left(\frac{x_0+\alpha y}{\|x_0+\alpha y\|}\right)\right]\Big|_{\alpha=0}=\frac{\dd}{\dd\alpha}\left[\frac{1}{1+\alpha^2}\langle x_0+\alpha y, A(x_0+\alpha y)\rangle\right]\Big|_{\alpha=0}=\\ 
     &=\frac{\dd}{\dd\alpha}\left[\frac{1}{1+\alpha^2}(a(x_0)+\alpha^2a(y)+ Re(\alpha\langle y, Ax_0\rangle))\right]\Big|_{\alpha=0}=\\ 
     &=\left[-\frac{2\alpha a(x_0)}{(1+\alpha^2)^2}+\frac{2xa(y)}{(1+x^2)^2}+Re\langle y, Ax_0\rangle \right]_{\alpha=0}=Re(\langle y, Ax_0\rangle)
  \end{align*}
  Zastępując $y$ przez $iy$ dostajemy $Im\langle y, Ax_0\rangle=0$, więc $\langle y, Ax_0\rangle=0$. Czyli $Ax_0\perp y$ dla każdego $y\perp x_0$, więc $Ax_0$ jest w przestrzeni rozpiętej przez $\{x_0\}$, czyli musi być wielokrotnością $x_0$. Stąd $x_0$ jest wektorem własnym $A$.
\end{proof}

\begin{theorem}{spektralne}{}
  Niech $A=A^*$ będzie operatorem Hilberta-Schmidta na przestrzeni $\Hh$. Istnieje wtedy baza ortonormalna $\{e_n\}$ przestrzeni Hilberta $\Hh$ oraz ciąg liczb $\{\lambda_n\}$ takich, że
  \begin{itemize}
    \item $Ax_n=\lambda_nx_n$, czyli $x_n$ jest wektorem własnym dla wartości własnej $\lambda_n$,
    \item jeśli $\dim(\Hh)=\infty$, to $\lim\lambda_n=0$.
  \end{itemize}
\end{theorem}
  W szczególności jeśli $\lambda\neq 0$ to $\{x\in\Hh\;:\;Ax=\lambda x\}$ jest skończenie wymiarową przestrzenią.

\begin{proof}
  Rozpatrujemy funkcję $a(x)$ jak wcześniej. Pokazaliśmy już w lemacie \ref{lmm:a jest ciagla}, że funkcja $a$ jest ciągła na kuli jednostkowej $B$ w słabej topologii. Z lematu \ref{lmm:kula jednostkowa slabo zwarta} wiemy, że $B$ jest zwarta. W takim razie $\sup_{x\in B} a(x)$ oraz $\inf_{y\in B}a(y)$ są osiągane przez pewne $x_0,y_0\in B$.
  
  Oznaczmy $M^+=\sup a(x)$, $M^-=\inf a(x)$. Mamy trzy opcje
  \begin{enumerate}
    \item $|M^+|\geq|M^-|$, wtedy korzystamy z lematu \ref{lmm:supremum to wartosc wlasna} tak jak stoi,
    \item $|M^+|\leq|M^-|$, wtedy musimy przeformułować lemat \ref{lmm:supremum to wartosc wlasna} dla $\inf$ (dowód jest analogiczny),
    \item $|M^+|=|M^-|=0$, wtedy operator jest zerowy i nie ma co dowodzić.
  \end{enumerate}
  W sytuacji 1. i 2. dostajemy wektor $x_0$ który osiąga pożądane przez nas ekstremum i jest wartością własną operatora $A$ dla $M^\pm$. Stawiamy go jako kolejny wektor z szukanej przez nas bazy ortonormalnej zawierającej wektory własne. 

  Ponieważ dopełnienie ortogonalne podprzestrzeni niezmienniczej jest również podprzestrzenią niezmienniczą, to $\{x_0\}^\perp\cap B$ jest nadal zwartą przestrzenią na której określony jest operator $A$. Co więcej, zbiór ten jest nadal słabo zwarty, więc możemy powtórzyć dla $B$ obciętego do $\{x_0\}^\perp\cap B$ powtórzyć szukanie wektora i wartości własnej. 
  
  Tą indukcyjną metodą dostajemy ciąg wektorów $\{x_n\}$ oraz ich wartości własnych $\{\lambda_n\}$, które spełniają $|\lambda_1|\geq|\lambda_2|\geq...>0$, co kończy dowód.
\end{proof}




